\SeventhArticle{5}

Da die Formel $\Delta(AR+\allowbreak A'R'+\allowbreak A'R'+\allowbreak A''R''+\allowbreak \text{usw.}+\allowbreak A^{(n)}R^{(n)})$ das Integral $\int y\,\mathrm{d}x$ von $x=g$ bis $x=g+\Delta$ oder das Integral $\Delta\int y\,\mathrm{d}t$ von $t=0$ bis $t=1$ genau darstellt, wenn $y$, in eine Reihe entwickelt, die Potenz $t^n$ nicht überschreitet, dagegen nur näherungsweise, wenn $y$ darüber hinausgeht, müssen wir noch zeigen, wie der durch die unmittelbar folgenden Glieder verursachte Fehler zu bestimmen ist. Wir bezeichnen allgemein mit $k^{(m)}$ die Differenz zwischen dem wahren Wert des Integrals $\int t^m\,\mathrm{d}t$ von $t=0$ bis $t=1$ und dem durch die Formel erhaltenen Wert, sodass
\[
\begin{aligned}
k&=1-R-R'-R''-R'''-\text{usw.}-R^{(n)}\\
k'&=\frac12-\frac1n(R'+2R''+3R'''+\text{usw.}+nR^{(n)})\\
k''&=\frac13-\frac1{nn}(R'+4R''+9R'''+\text{usw.}+nnR^{(n)})\\
k'''&=\frac14-\frac1{n^2}(R'+8R''+27R'''+\text{usw.}\ n^3R^{(n)})
\end{aligned}
\]
usw. Wird also $y$ in die Reihe
\[
K+K't+K''tt+K'''t^3+\text{usw.}
\]
entwickelt, so stellt der folgende Ausdruck offenbar die Differenz zwischen dem wahren Wert des Integrals $\int y\,\mathrm{d}t$ und dem durch die Formel gelieferten Näherungswert dar:
\[
Kk+K'k'+K''k''+K'''k'''+\text{usw.}
\]
Aber $k$, $k'$, $k''$ usw. bis $k^{(n)}$ verschwinden offenbar von selbst ($=0$); die Korrektur der Näherungsformel lautet daher
\[
K^{(n+1)}k^{(n+1)}+K^{(n+2)}k^{(n+2)}+K^{(n+3)}k^{(n+3)}+\text{usw.}
\]

Die Beschaffenheit der Größen $k^{(n+1}$, $k^{(n+2)}$ usw. werden wir unten genauer untersuchen. Hier genügt es, für die einzelnen Werte von $n$ die Zahlenwerte der ersten oder der zweiten anzugeben, damit sich daraus der Genauigkeitsgrad der Näherungsformel beurteilen lässt.
\SourceLeaf{179}

\begin{longtable}{@{}l@{\quad}l@{}}
$\text{Für }n=1\ \text{haben wir }$ & $k''=-\frac16,\quad k'''=-\frac14,\quad k''''=-\frac3{10}$ \\[16pt]
$\text{Für }n=2\ \text{finden wir }$ & $k'''=0,\quad k''''=-\frac1{120},\quad k^{\mathrm{V}}=-\frac1{48}$ \\[16pt]
$\text{Für }n=3\ \text{ergibt sich }$ & $k''''=-\frac1{270},\quad k^{\mathrm{V}}=-\frac1{108}$ \\[16pt]
$\text{Für }n=4\ldotp\ldotp\ldotp\ $ & $k^{\mathrm{V}}=0,\quad k^{\mathrm{VI}}=-\frac1{2688},\quad k^{\mathrm{VII}}=-\frac1{768}$ \\[16pt]
$\text{Für }n=5\ldotp\ldotp\ldotp\ $ & $k^{\mathrm{VI}}=-\frac{11}{52500},\quad k^{\mathrm{VII}}=-\frac{11}{15000}$ \\[16pt]
$\text{Für }n=6\ldotp\ldotp\ldotp\ $ & $k^{\mathrm{VII}}=0,\quad k^{\mathrm{VIII}}=-\frac1{38880},\quad k^{\mathrm{IX}}=-\frac1{8640}$ \\[16pt]
$\text{Für }n=7\ldotp\ldotp\ldotp\ $ & $k^{\mathrm{VIII}}=-\frac{167}{10588410},\quad k^{\mathrm{IX}}=-\frac{167}{2352980}$ \\[16pt]
$\text{Für }n=8\ldotp\ldotp\ldotp\ $ & $\begin{aligned}[t]&k^{\mathrm{IX}}=0,\quad k^{\mathrm{X}}=-\frac{37}{17301504},\\[3pt]&k^{\mathrm{XI}}=-\frac{37}{3145728}\end{aligned}$ \\[16pt]\noalign{\vskip 8pt}
$\text{Für }n=9\ldotp\ldotp\ldotp\ $ & $\begin{aligned}[t]&k^{\mathrm{X}}=-\frac{865}{631351908},\\[3pt]&k^{\mathrm{XI}}=-\frac{865}{114791256}\end{aligned}$ \\[16pt]\noalign{\vskip 8pt}
$\text{Für }n=10\ldotp\ldotp\ldotp\ $ & $\begin{aligned}[t]&k^{\mathrm{XI}}=0,\quad k^{\mathrm{XII}}=-\frac{26927}{136500000000},\\[3pt]&k^{\mathrm{XIII}}=-\frac{26927}{21000000000}\end{aligned}$ \\[16pt]
\end{longtable}




Bei einem geraden Wert von $n$ beobachten wir hier überall $k^{(n+1)}=0$ und außerdem $k^{(n+3)}=\frac{n+3}{2}k^{(n+2)}$; bei einem ungeraden Wert von $n$ ergibt sich dagegen überall $k^{(n+2)}=\frac{n+2}{2}k^{(n+1)}$. Der Grund dieser Ergebnisse lässt sich leicht aus den folgenden Betrachtungen entnehmen.

Wir bezeichnen allgemein mit $l^{(m)}$ die Differenz zwischen dem wahren Wert des Integrals $\int(t-\frac12)^m\,\mathrm{d}t$ von $t=0$ bis $t=1$ und dem durch die Näherungsformel erhaltenen Wert, sodass
\[
\begin{aligned}
l^{(m)}={}&\int(t-\tfrac12)^m\,\mathrm{d}t
-\bigl[(-\tfrac12)^mR+(\tfrac1n-\tfrac12)^mR'
\\ &\qquad +(\tfrac2n-\tfrac12)^mR''+(\tfrac3n-\tfrac12)^mR'''+\text{usw.}\\
&\qquad +(\tfrac12-\tfrac1n)^mR^{(n-1)}+(\tfrac12)^mR^{(n)}\bigr]
\end{aligned}
\]
wobei das Integral von $t=0$ bis $t=1$ genommen wird. Offenbar verschwinden bei einem ungeraden Wert von $m$ sowohl der wahre als auch der angenäherte Wert des Integrals. Es ist daher $l'=0$, $l'''=0$, $l^{\mathrm{V}}=0$, $l^{\mathrm{VII}}=0$ usw., oder allgemein $l^{(m)}=0$ bei einem ungeraden Wert von $m$. Bei einem geraden Wert von $m$ geben wir der Formel dagegen die folgende Gestalt
\[
\begin{aligned}
l^{(m)}={}&\frac1{2^m(m+1)}-\frac2{n^m}\bigl((\tfrac12n)^mR
+(\tfrac12n-1)^mR'\\ &\qquad +(\tfrac12n-2)^mR''+\text{usw.}\\
&\qquad +2^mR^{(\frac12n-2)}+R^{(\frac12n-1)}\bigr)
\end{aligned}
\]
wenn zugleich $n$ gerade ist; oder die folgende Gestalt
\[
\begin{aligned}
l^{(m)}={}&\frac1{2^m}\bigl(\frac1{m+1}-\frac2{n^m}\bigl(n^mR
+(n-2)^mR'\\ &\qquad +(n-4)^mR''+\text{usw.}\\
&\qquad +3^mR^{(\frac12n-\frac32)}+R^{(\frac12n-\frac12)}\bigr)\bigr)
\end{aligned}
\]
wenn zugleich $n$ ungerade ist.
\par\SourceLeaf{180}
Ergibt sich also bei der Entwicklung von $y$ in eine Reihe nach Potenzen von $t-\frac12$
\[
y=L+L'(t-\tfrac12)+L''(t-\tfrac12)^2+L'''(t-\tfrac12)^3+\text{usw.}
\]
so lautet die am Näherungswert des Integrals $\int y\,\mathrm{d}t$ von $t=0$ bis $t=1$ anzubringende Korrektur
\[
Ll+L''l''+L''''l''''+L^{\mathrm{VI}}l^{\mathrm{VI}}+\text{usw.}
\]
\Needspace{14\baselineskip}oder vielmehr, da $l^{(m)}$ für jeden ganzzahligen Wert von $m$, der nicht größer als $n$ ist, notwendig verschwindet, lautet die Korrektur
\[
L^{(n+2)}l^{(n+2)}+L^{(n+4)}l^{(n+4)}+L^{(n+6)}l^{(n+6)}+\text{usw.}
\]
bei geradem $n$ oder
\[
L^{(n+1)}l^{(n+1)}+L^{(n+3)}l^{(n+3)}+L^{(n+5)}l^{(n+5)}+\text{usw.}
\]
bei ungeradem $n$.

Die Korrekturen $l^{(m)}$ lassen sich nun sehr leicht auf $k^{(m)}$ zurückführen und umgekehrt. Da nämlich
\[
(t-\tfrac12)^m=t^m-\tfrac12m\mathbin{.}t^{m-1}
+\tfrac14\mathbin{.}\frac{m(m-1)}{1\mathbin{.}2}t^{m-2}+\text{usw.}
\]
so ergibt sich
\[
l^{(m)}=k^{(m)}-\tfrac12m k^{(m-1)}
+\tfrac14\mathbin{.}\frac{m(m-1)}{1\mathbin{.}2}k^{(m-2)}+\text{usw.}
\]
Ebenso ergibt sich
\[
k^{(m)}=l^{(m)}+\tfrac12m l^{(m-1)}
+\tfrac14\mathbin{.}\frac{m(m-1)}{1\mathbin{.}2}l^{(m-1)}+\text{usw.}
\]



Aus der letzteren Formel werden die Glieder entfernt, in denen $l$ einen ungeraden Index hat; beide Formeln sind jedoch nur bis zum Index $n+1$ einschließlich fortzusetzen. Wir haben daher offenbar
\[
\begin{aligned}
&\textit{bei geradem $n$}\\
&\qquad k^{(n+1)}=0\\
&\qquad k^{(n+2)}=l^{(n+2)}\\
&\qquad k^{(n+3)}=\frac{n+3}{2}\mathbin{.}l^{(n+2)}\\
&\textit{bei ungeradem $n$}\\
&\qquad k^{(n+1)}=l^{(n+1)}\\
&\qquad k^{(n+2)}=\frac{n+2}{2}\mathbin{.}l^{(n+1)}
\end{aligned}
\]
woraus die oben angegebenen Beobachtungen folgen.
\SourceLeaf{181}
\SeventhArticle{6}

Im Allgemeinen wird es daher bei der Anwendung des Verfahrens von Cotes vorzuziehen sein, $n$ einen geraden Wert zu geben, also eine ungerade Anzahl von Gliedern zu verwenden. Die Genauigkeit nimmt nämlich nur sehr wenig zu, wenn wir von einem geraden Wert von $n$ zum nächsthöheren ungeraden übergehen, da der Fehler von derselben Ordnung bleibt, wenn auch sein Koeffizient etwas kleiner ist. Geht man dagegen von einem ungeraden Wert von $n$ zum nächstfolgenden geraden über, so erhöht sich die Ordnung des Fehlers um zwei, und ein stärker verminderter Koeffizient erhöht die Genauigkeit zusätzlich. Werden beispielsweise fünf Glieder verwendet, also bei $n=4$, so wird der Fehler näherungsweise durch $-\frac{1}{2688}K^6$ oder durch $-\frac{1}{2688}L^6$ ausgedrückt. Setzt man $n=5$, so ist der Fehler näherungsweise $-\frac{11}{52500}K^6$ oder $-\frac{11}{52500}L^6$ und somit noch nicht einmal auf die Hälfte des vorigen Fehlers herabgesetzt. Setzt man dagegen $n=6$, so ergibt sich für den Fehler näherungsweise $=-\frac{1}{38880}K^8$ oder $=-\frac{1}{38880}L^8$; die Genauigkeit nimmt dabei umso mehr zu, je schneller die Reihe, in die die Funktion entwickelt wird, bereits von sich aus konvergiert.

\SeventhArticle{7}

Nach diesen vorausgeschickten Bemerkungen über das Verfahren von Cotes gehen wir zu einer allgemeinen Untersuchung über und lassen die Bedingung fallen, dass die Werte von $x$ eine arithmetische Folge bilden. Wir behandeln also die Aufgabe, den Wert des Integrals $\int y\,\mathrm{d}x$ zwischen gegebenen Grenzen aus einigen gegebenen Werten von $y$ entweder genau oder möglichst nahe zu bestimmen. Wir nehmen an, das Integral sei von $x=g$ bis $x=g+\Delta$ zu nehmen, und führen anstelle von $x$ eine andere Variable $t=\frac{x-g}{\Delta}$ ein, sodass das Integral $\Delta\int y\,\mathrm{d}t$ von $t=0$ bis $t=1$ zu untersuchen ist. Die $n+1$ gegebenen Werte von $y$, nämlich $A,\ A',\ A'',\ A'''\ldotp\ldotp\ldotp\ldotp A^{(n)}$, sollen den verschiedenen Werten von $t$, nämlich $a,\ a',\ a'',\ a'''\ldotp\ldotp\ldotp\ldotp a^{(n)}$, entsprechen, und wir bezeichnen mit $Y$ das folgende Polynom vom Grad $n$:
\[
\begin{aligned}
& A\;\frac{(t-a')(t-a'')(t-a''')\ldotp\ldotp\ldotp\ldotp(t-a^{(n)})}
{(a-a')(a-a'')(a-a''')\ldotp\ldotp\ldotp\ldotp(a-a^{(n)})}\\[6pt]
&{}+A'\;\frac{(t-a)(t-a'')(t-a''')\ldotp\ldotp\ldotp\ldotp(t-a^{(n)})}
{(a'-a)(a'-a'')(a'-a''')\ldotp\ldotp\ldotp\ldotp(a'-a^{(n)})}\\[6pt]
&{}+A''\;\frac{(t-a)(t-a')(t-a''')\ldotp\ldotp\ldotp\ldotp(t-a^{(n)})}
{(a''-a)(a''-a')(a''-a''')\ldotp\ldotp\ldotp\ldotp(a''-a^{(n)})}\\[6pt]
&{}+\text{usw.}\\[6pt]
&{}+A^{(n)}\;\frac{(t-a)(t-a')(t-a'')\ldotp\ldotp\ldotp\ldotp(t-a^{(n-1)})}
{(a^{(n)}-a)(a^{(n)}-a')(a^{(n)}-a'')\ldotp\ldotp\ldotp\ldotp(a^{(n)}-a^{(n-1)})}
\end{aligned}
\]
Offenbar stimmen die Werte dieser Funktion, wenn $t$ einer der Größen $a,\allowbreak \ a',\allowbreak \ a'',\allowbreak \ a'''\ldotp\ldotp\ldotp\ldotp a^{(m)}$ gleichgesetzt wird, mit den entsprechenden Werten der Funktion $y$ überein. Daraus\SourceLeaf{182} schließen wir wie in Artikel 2, dass $Y$ mit $y$ identisch ist, wenn $y$ ebenfalls ein Polynom ist, dessen Grad $n$ nicht überschreitet, oder dass sie zumindest an die Stelle von $y$ treten kann, wenn $y$, in eine Reihe nach Potenzen von $t$ entwickelt, so schnell konvergiert, dass Glieder höherer Ordnungen vernachlässigt werden dürfen.


